\meetingheader
    {Mar 09, 2026}
    {}
    {\meetinglinks{
        \meetinglink{Smelly Tool Descriptions (Feb 2026)}{https://arxiv.org/abs/2602.14878}
    }}
    {\begin{itemize}
        \item Pipeline fixes: dropped DeepSpeed for HF Trainer, fixed image padding, replaced GPT-4o judge with local Qwen3.5-27B-FP8, fixed batch size OOM, made training limits configurable
        \item Tool description design: restructured 20+ tools to Purpose/Guidelines/Limitations per smelly-descriptions ablation
        \item 5 agent episodes on FineVision5: Ep~1 matches baseline (MMStar 30.3, OCRBench 333); failure modes from aggressive filtering and budget reduction
    \end{itemize}}
    {}



%% ============================================================
%% SECTION: TOOL DESCRIPTION DESIGN
%% ============================================================

\newpage
\section{Tool Description Design}

\subsection{How tool definitions reach the model}

In tool-calling LLM systems, tool definitions (name, description, parameter schema) are not passed through a side channel. They are serialized into the model's context window as literal prompt text, placed in either the system or user turn depending on the provider\footnote{See e.g.\ the OpenAI Harmony format specification: \url{https://developers.openai.com/cookbook/articles/openai-harmony}} (Figure~\ref{fig:tool-formats}). HuggingFace Transformers exposes this via the \texttt{tools} parameter of \texttt{apply\_chat\_template}, which passes tool JSON schemas through the model's Jinja2 template to produce the final prompt string.

This means tool description quality has a direct effect on agent behavior: a vague or noisy description wastes context tokens and increases the chance of incorrect tool selection or malformed arguments. The description is, in effect, the API contract between the tool developer and the LLM policy.

% Sources: HF chat_templating docs (https://huggingface.co/docs/transformers/en/chat_templating_writing),
%   HF tool use docs (https://huggingface.co/docs/transformers/main/en/chat_extras),
%   OpenAI Harmony format (https://developers.openai.com/cookbook/articles/openai-harmony),
%   Meta Llama 3.1 prompt format (https://www.llama.com/docs/model-cards-and-prompt-formats/llama3_1/)

\definecolor{anthropicbg}{HTML}{FFF0E6}
\definecolor{anthropicframe}{HTML}{D97706}
\definecolor{harmonybg}{HTML}{E8F5E9}
\definecolor{harmonyframe}{HTML}{388E3C}
\definecolor{qwenbg}{HTML}{E8F1FC}
\definecolor{qwenframe}{HTML}{2A74B9}
\definecolor{llamabg}{HTML}{F3E8FF}
\definecolor{llamaframe}{HTML}{7C3AED}

\begin{figure}[H]
\centering
\small
\begin{minipage}{0.8\linewidth}

\begin{tcolorbox}[enhanced, colback=anthropicbg, colframe=anthropicframe, boxrule=0.8pt, left=4pt, right=4pt, top=4pt, bottom=3pt, sharp corners,
  attach boxed title to top left={xshift=4pt, yshift=-2pt},
  boxed title style={empty, size=minimal},
  coltitle=anthropicframe, fonttitle=\sffamily\bfseries\large,
  title={Anthropic (Claude)}]
\ttfamily\scriptsize
<functions>\\
<function>\{"name": "vlm\_filter",\\
\quad "description": \colorbox{yellow!30}{"Remove rows whose images}\\
\quad \colorbox{yellow!30}{are corrupt, too small, or extreme aspect}\\
\quad \colorbox{yellow!30}{ratios."},\\
\quad "parameters": \{"dataset\_name":\\
\quad\quad \{"type": "string"\}\}\}</function>\\
</functions>
\end{tcolorbox}

\vspace{-4pt}

\begin{tcolorbox}[enhanced, colback=harmonybg, colframe=harmonyframe, boxrule=0.8pt, left=4pt, right=4pt, top=4pt, bottom=3pt, sharp corners,
  attach boxed title to top left={xshift=4pt, yshift=-2pt},
  boxed title style={empty, size=minimal},
  coltitle=harmonyframe, fonttitle=\sffamily\bfseries\large,
  title={gpt-oss-120b}]
\ttfamily\scriptsize
<|start|>developer<|message|>\\
\# Tools\\
\\
\#\# functions\\
\\
namespace functions \{\\
\\
// \colorbox{yellow!30}{Remove rows whose images are corrupt, too small,}\\
// \colorbox{yellow!30}{or extreme aspect ratios.}\\
type vlm\_filter = (\_: \{\\
dataset\_name: string,\\
\}) => any;\\
\\
\} // namespace functions<|end|>
\end{tcolorbox}

\vspace{-4pt}

\begin{tcolorbox}[enhanced, colback=qwenbg, colframe=qwenframe, boxrule=0.8pt, left=4pt, right=4pt, top=4pt, bottom=3pt, sharp corners,
  attach boxed title to top left={xshift=4pt, yshift=-2pt},
  boxed title style={empty, size=minimal},
  coltitle=qwenframe, fonttitle=\sffamily\bfseries\large,
  title={Qwen 3.5}]
\ttfamily\scriptsize
<|im\_start|>system\\
\# Tools\\
\\
You have access to the following functions:\\
\\
<tools>\\
\{"type": "function", "function": \{"name": "vlm\_filter",\\
"description": \colorbox{yellow!30}{"Remove rows whose images are corrupt,}\\
\colorbox{yellow!30}{too small, or extreme aspect ratios."},\\
"parameters": \{"dataset\_name": \{"type": "string"\}\}\}\}\\
</tools><|im\_end|>
\end{tcolorbox}

\vspace{-4pt}

\begin{tcolorbox}[enhanced, colback=llamabg, colframe=llamaframe, boxrule=0.8pt, left=4pt, right=4pt, top=4pt, bottom=3pt, sharp corners,
  attach boxed title to top left={xshift=4pt, yshift=-2pt},
  boxed title style={empty, size=minimal},
  coltitle=llamaframe, fonttitle=\sffamily\bfseries\large,
  title={Llama 3.1}]
\ttfamily\scriptsize
<|start\_header\_id|>user<|end\_header\_id|>\\
\\
Given the following functions, please respond with\\
a JSON for a function call ...\\
\\
\{"type": "function", "function": \{"name": "vlm\_filter",\\
"description": \colorbox{yellow!30}{"Remove rows whose images are corrupt,}\\
\colorbox{yellow!30}{too small, or extreme aspect ratios."},\\
"parameters": \{"dataset\_name": \{"type": "string"\}\}\}\}\\
\\
Filter the dataset<|eot\_id|>
\end{tcolorbox}

\end{minipage}
\caption{One tool schema, four prompt formats. The highlighted description appears verbatim in each. Tool descriptions are literal prompt text, not metadata.}
\label{fig:tool-formats}
\end{figure}

\subsection{Restructuring tool descriptions}

\definecolor{purposebg}{HTML}{FDECEA}
\definecolor{guidelinesbg}{HTML}{FFF9E6}
\definecolor{limitationsbg}{HTML}{E6F0FA}

Hasan et al.~\cite{smellyToolDesc2026} synthesize six tool description components (Purpose, Guidelines, Limitations, Parameter Explanation, Examples, Length \& Completeness) from Anthropic's official docs, practitioner blogs, and prior work, then ablate them across 856 MCP tools and four models. Across four models and six domains (finance, 3D design, browser automation, location navigation, repo management, web search), agents with all six components achieve +5.85 pp median success rate and +15.12\% average evaluator score, at the cost of 67.46\% more execution steps. However, 16.67\% of domain-model combinations regress, Purpose is essential in every effective combination, and the best combination is domain- and model-dependent. Their strongest reduced set, P+G+L+PEx (everything minus Examples), was stable in 60\% of domain-model pairs. I dropped Parameter Explanation as well because our tools have simple parameters (typically just a \texttt{dataset\_id}), leaving three components: \colorbox{purposebg}{\textbf{Purpose}} (what the tool does and returns), \colorbox{guidelinesbg}{\textbf{Guidelines}} (when to use it and ordering constraints), and \colorbox{limitationsbg}{\textbf{Limitations}} (failure modes and performance constraints). Below are three verbatim docstrings from the codebase showing this format applied to tools at different pipeline stages.

\medskip
\noindent\textbf{\texttt{vlm\_filter}} \textit{(preselection, image structural quality)}
\begin{tcolorbox}[enhanced, colback=white, colframe=black!60, boxrule=0.6pt, left=6pt, right=6pt, top=0pt, bottom=0pt, sharp corners, boxsep=0pt]
\colorbox{purposebg}{\parbox{\dimexpr\linewidth-2\fboxsep}{%
\ttfamily\scriptsize\vspace{4pt}
\textbf{Purpose}: Remove rows whose images are corrupt, too
small (any side <200px), or have extreme aspect ratios
(>3:1). Returns kept/rejected counts and per-reason
breakdown.\vspace{4pt}
}}\\[-1pt]
\colorbox{guidelinesbg}{\parbox{\dimexpr\linewidth-2\fboxsep}{%
\ttfamily\scriptsize\vspace{4pt}
\textbf{Guidelines}: Run on each dataset after profile\_datasets
and before text-based filters. Only requires the
dataset\_id; thresholds are fixed. Safe on all
dataset types.\vspace{4pt}
}}\\[-1pt]
\colorbox{limitationsbg}{\parbox{\dimexpr\linewidth-2\fboxsep}{%
\ttfamily\scriptsize\vspace{4pt}
\textbf{Limitations}: Requires loading full image data into
memory (calls ensure\_full internally). Slow on large
datasets with many shards. Does not assess image
content quality, only structural properties.\vspace{4pt}
}}
\end{tcolorbox}

\medskip
\noindent\textbf{\texttt{submit\_finetune}} \textit{(postselection, training)}
\begin{tcolorbox}[enhanced, colback=white, colframe=black!60, boxrule=0.6pt, left=6pt, right=6pt, top=0pt, bottom=0pt, sharp corners, boxsep=0pt]
\colorbox{purposebg}{\parbox{\dimexpr\linewidth-2\fboxsep}{%
\ttfamily\scriptsize\vspace{4pt}
\textbf{Purpose}: Finetune the target model on the curated dataset. Blocks until training completes. Returns a job\_id (e.g.\ ``ft-abc123'') that identifies the trained model for evaluation.\vspace{4pt}
}}\\[-1pt]
\colorbox{guidelinesbg}{\parbox{\dimexpr\linewidth-2\fboxsep}{%
\ttfamily\scriptsize\vspace{4pt}
\textbf{Guidelines}: Run after save\_to\_disk. Pass the dataset\_id of the mixed dataset (e.g.\ ``vqav2+captcha+cosyn\_chart''). All training hyperparameters (learning rate, epochs, batch size) are fixed by the task YAML. You only control which data to train on. The returned job\_id is required by submit\_eval.\vspace{4pt}
}}\\[-1pt]
\colorbox{limitationsbg}{\parbox{\dimexpr\linewidth-2\fboxsep}{%
\ttfamily\scriptsize\vspace{4pt}
\textbf{Limitations}: Blocking call. Typically takes 10--30 minutes depending on dataset size. Training is capped at max\_steps from the task config. Cannot be cancelled once started. All hyperparameters are fixed. The agent has no control over training configuration.\vspace{4pt}
}}
\end{tcolorbox}

\medskip
\noindent\textbf{\texttt{submit\_eval}} \textit{(postselection, evaluation)}
\begin{tcolorbox}[enhanced, colback=white, colframe=black!60, boxrule=0.6pt, left=6pt, right=6pt, top=0pt, bottom=0pt, sharp corners, boxsep=0pt]
\colorbox{purposebg}{\parbox{\dimexpr\linewidth-2\fboxsep}{%
\ttfamily\scriptsize\vspace{4pt}
\textbf{Purpose}: Evaluate a finetuned model on VLM benchmarks configured in the task YAML (e.g.\ MMStar, OCRBench, MMMU, HallusionBench, MMVet, MathVista). Blocks until all benchmarks complete. Returns per-benchmark scores normalized to [0,1].\vspace{4pt}
}}\\[-1pt]
\colorbox{guidelinesbg}{\parbox{\dimexpr\linewidth-2\fboxsep}{%
\ttfamily\scriptsize\vspace{4pt}
\textbf{Guidelines}: Run after submit\_finetune. Pass model\_id set to the job\_id returned by submit\_finetune (e.g.\ ``ft-abc123''). Do NOT pass a filesystem path. Benchmarks are configured automatically from the task YAML. Use the returned scores to inform the next episode's data curation strategy.\vspace{4pt}
}}\\[-1pt]
\colorbox{limitationsbg}{\parbox{\dimexpr\linewidth-2\fboxsep}{%
\ttfamily\scriptsize\vspace{4pt}
\textbf{Limitations}: Blocking call. Evaluation time scales with number of benchmarks (typically 15--45 minutes total). Benchmarks and evaluation protocol are fixed by the task config. The agent cannot select which benchmarks to run. Requires a running vLLM inference server and judge model.\vspace{4pt}
}}
\end{tcolorbox}


%% ============================================================
%% SECTION 2: EXPERIMENTAL SETUP
%% ============================================================

\newpage
\section{Experimental Setup}
\label{sec:setup}

\subsection{Data and fine-tuning}

Five datasets from FineVision5, each containing 10k samples of image--text pairs (VQA, captioning, OCR, instruction-following). The agent's sample budget is 10k total, so it must select and mix from the available 50k. One dataset (captcha) is consistently auto-excluded by \texttt{vlm\_filter} due to structural image failures, leaving 4 effective source datasets.

The target model is LLaVA-1.5-7B, full fine-tuned on a single H200 (140GB). Vision tower frozen; only the language model and multimodal projector are trained.

Two fixes from the previous setup: replaced DeepSpeed with plain HF Trainer + Liger kernel~\cite{hsu2024liger} for simpler single-GPU training, and restored LLaVA's \texttt{expand2square()} image padding to prevent center-crop from losing edge content.

\begin{table}[H]
\centering
\small
\caption{Fine-tuning hyperparameters. All fixed by the task configuration. The agent controls only which data is selected.}
\label{tab:finetune-config}
\begin{tabular}{ll}
\toprule
\textbf{Parameter} & \textbf{Value} \\
\midrule
Epochs & 1 \\
Batch size (per device) & 2 \\
Gradient accumulation & 2 (effective batch 4) \\
Learning rate & 2e-5 (cosine, warmup 3\%) \\
Weight decay & 0.0 \\
Precision & bf16 \\
Optimizer & AdamW (fused) \\
Gradient checkpointing & enabled \\
Liger kernel & fused RoPE, RMSNorm, SwiGLU, CE \\
Max steps & 1{,}250 (cap), timeout 1{,}800s \\
Max sequence length & 2{,}048 tokens \\
Step time & $\sim$1.0s/step \\
\bottomrule
\end{tabular}
\end{table}

\subsection{Evaluation}

Six benchmarks via VLMEvalKit, with Qwen3.5-27B-FP8 as the LLM judge (served locally via vLLM). Inference uses vLLM for the fine-tuned model.

\begin{table}[H]
\centering
\small
\caption{LLaVA-1.5-7B zero-shot baselines. Published = from original papers; Ours = VLMEvalKit with Qwen3.5-27B-FP8 judge, vLLM inference. Differences are within normal eval-toolkit variance.}
\label{tab:judge-sanity}
\begin{tabular}{lccll}
\toprule
\textbf{Benchmark} & \textbf{Published} & \textbf{Ours} & \textbf{Source} & \textbf{Note} \\
\midrule
MMStar & 30.3 & 33.3 & \cite{chen2024mmstar} & paper Table~5 \\
OCRBench & 297 & 335 & \cite{liu2023ocrbench} & paper; ours higher due to VLMEvalKit prompt diffs \\
HallusionBench & --- & 25.2 & --- & composite $(\text{aAcc}+\text{fAcc}+\text{qAcc})/3$; no paper value for 7B \\
MMMU\_DEV\_VAL & 36.4 & 36.2 & \cite{yue2024mmmu} & paper (lmms-eval reports 35.3) \\
MMVet & 31.1 & --- & \cite{liu2024llava15} & LLaVA-1.5 paper Table~4; not in our baseline run \\
MathVista\_MINI & --- & --- & --- & original LLaVA-13B scores 26.1~\cite{lu2024mathvista}; no paper value for 1.5-7B \\
\bottomrule
\end{tabular}
\end{table}

\noindent The reward signal used by the agent is a normalized composite of benchmark scores. Judge-dependent benchmarks (MMVet, MathVista) use Qwen3.5-27B-FP8 served locally via vLLM.


%% ============================================================
%% SECTION 3: EPISODE RESULTS
%% ============================================================

\newpage
\section{Episode Results}

With all fixes applied, the agent runs the full loop on the setup described in Section~\ref{sec:setup}.


\paragraph{Episode summary.}

\begin{table}[H]
\centering
\small
\caption{Episode behavior summary. Steps = total tool calls. Pipeline shows the tool sequence. Ep~5 is the first episode with restructured tool descriptions and seed memories from Ep~1--4.}
\label{tab:episode-summary}
\begin{tabular}{clcl}
\toprule
\textbf{Ep} & \textbf{Pipeline} & \textbf{Steps} & \textbf{Outcome} \\
\midrule
1 & profile $\to$ explore(4) $\to$ vlm\_filter(4) $\to$ quality\_filter(4) $\to$ mix $\to$ ft $\to$ eval & 22 & done (reward 0.293) \\
2 & profile $\to$ explore(5) $\to$ vlm\_filter(5) $\to$ quality\_filter(5) $\to$ mix $\to$ ft & 20 & eval crashed \\
3 & profile $\to$ vlm\_filter(5) $\to$ quality\_filter(4) $\to$ rollback(2) $\to$ mix $\to$ ft(2) $\to$ eval & 25 & done (reward 0.204) \\
4 & profile $\to$ vlm\_filter(5) $\to$ mix(10k) $\to$ ft(timeout) $\to$ mix(5k) $\to$ ft $\to$ eval & 18 & eval crashed \\
\midrule
5 & profile $\to$ vlm\_filter(5) $\to$ mix $\to$ save $\to$ ft $\to$ eval & \textbf{12} & eval OOM (scores recovered) \\
\bottomrule
\end{tabular}
\end{table}

% Color-coded transcript: blue = agent reasoning, orange = tool execution
\definecolor{agentbg}{HTML}{E8F1FC}
\definecolor{agentframe}{HTML}{2563EB}
\definecolor{actionbg}{HTML}{FFF3E0}
\definecolor{actionframe}{HTML}{E65100}

\newtcolorbox{agentthought}[1][]{enhanced, before={\noindent}, colback=agentbg, colframe=agentframe, boxrule=0.6pt, left=6pt, right=6pt, top=4pt, bottom=3pt, sharp corners, fonttitle=\sffamily\bfseries\small, coltitle=agentframe, attach boxed title to top left={xshift=4pt, yshift=-2pt}, boxed title style={empty, size=minimal}, #1}

\newtcolorbox{toolaction}[1][]{enhanced, colback=actionbg, colframe=actionframe, boxrule=0.6pt, left=6pt, right=6pt, top=4pt, bottom=3pt, sharp corners, fonttitle=\sffamily\bfseries\small, coltitle=actionframe, attach boxed title to top left={xshift=4pt, yshift=-2pt}, boxed title style={empty, size=minimal}, #1}

\paragraph{Ep~5: verbatim reasoning trace.}
Episode~5 is the first run after the tool description rewrite (Section~1.2) and with seed memories from Ep~1--4. It completed the full pipeline in \textbf{12 steps}. All quotes below are verbatim from \texttt{events.jsonl} (167 lines) with line numbers cited. \colorbox{agentbg}{\textcolor{agentframe}{\textbf{Blue}}} = agent reasoning; \colorbox{actionbg}{\textcolor{actionframe}{\textbf{Orange}}} = tool execution.

\subparagraph{Step 0:Strategic planning (line 3).}
Before any tool call, the agent reviews all prior episodes from seed memory:

\begin{agentthought}[title={Agent Thought}]
\small\ttfamily
\textbf{Review of past episodes}\\
- Episode 1 (balanced 25\% each) achieved the highest MMStar score (\textasciitilde30\%) -- essentially the baseline. OCRBench was also decent (33.3). This suggests a uniform mix provides good general reasoning ability.\\
- Episodes 2-4 suffered from crashes or severe MMStar drops. Key issues observed:\\
\quad 1. \textbf{Over-aggressive quality\_filter} (especially in Episode 3) likely removed useful reasoning instructions, collapsing MMStar to \textasciitilde1\%.\\
\quad 2. \textbf{Dataset size reductions} (Episode 4 halved to 5k) dramatically harmed MMStar, confirming we need the full 10k budget.\\
\quad 3. \textbf{Zero usage of captcha} -- it was never sampled, yet OCRBench may benefit from OCR-focused data.\\[6pt]
\textbf{Proposed revised strategy}\\
2. \textbf{Image filtering only}: Apply vlm\_filter to each dataset to drop corrupt/too-small images. \textbf{Do NOT run quality\_filter} (or keep all thresholds at 0) to avoid unintentionally discarding valuable instruction text.\\
3. \textbf{Mix ratio computation}: Manually verify that captcha gets a non-zero share (e.g., \textasciitilde10-15\%) to help OCRBench, while keeping strong instruction sources (llava\_instruct, cosyn\_chart) at \textasciitilde25-30\% each.
\end{agentthought}

\noindent Key decision: the agent explicitly \textbf{skips text-based filtering}. It attributes Ep~3's MMStar collapse to \texttt{quality\_filter}, though the actual trace shows the agent rolled back aggressive filters and recovered the data --- the true cause of the 1\% score remains unclear (possibly the zombie-crashed first finetune or a \texttt{save\_to\_disk} row duplication bug).

\subparagraph{Steps 1--6:Profile and filter (lines 45--106).}
\mbox{}

\begin{toolaction}[title={Tool Calls}]
\scriptsize
\begin{tabular}{@{}rlrr@{}}
\toprule
\textbf{Step} & \textbf{Call} & \textbf{Kept/Total} & \textbf{Time} \\
\midrule
1 & \texttt{profile\_datasets(sample\_n=1000, seed=42)} & & 0.8s \\
2 & \texttt{vlm\_filter(vqav2)} & 9{,}970 / 10{,}000 & 20s \\
3 & \texttt{vlm\_filter(captcha)} & \textbf{0 / 10{,}000} & 4s \\
4 & \texttt{vlm\_filter(cosyn\_chart)} & 9{,}942 / 10{,}000 & 339s \\
5 & \texttt{vlm\_filter(llava\_instruct)} & 9{,}980 / 10{,}000 & 147s \\
6 & \texttt{vlm\_filter(ocrvqa)} & 9{,}829 / 10{,}000 & 14s \\
\bottomrule
\end{tabular}
\end{toolaction}

\noindent Captcha returns 0 rows. The agent's plan called for 10--15\% captcha, but it adapts immediately (line 71):

\begin{agentthought}[title={Agent Thought:data-driven plan revision}]
\small\ttfamily
Whoa! captcha all images filtered out (likely low-res). That means captcha may not be usable for visual instruction (maybe it's text captcha images small). We may still consider using it? But images are all too small; we cannot use them. So we should drop captcha from mix.
\end{agentthought}

\subparagraph{Step 7:Proportional reasoning (line 107).}
The agent filtered 10k-row samples but needs full-dataset effective sizes for mixing. It derives proportional estimates unprompted:

\begin{agentthought}[title={Agent Thought:quantitative reasoning}]
\small\ttfamily
We don't have full counts after filtering, but we know original dataset sizes from profile (num\_rows). We can estimate filtered proportion = kept/original\_rows * total rows.\\[4pt]
From profile: vqav2 total 82772 rows, kept proportion = 9970/10000=0.997, so filtered rows \ensuremath{\approx} 82772*0.997 \ensuremath{\approx} 82540.\\
captcha total 113062 rows, kept proportion =0, filtered rows =0.\\
cosyn\_chart total 116814 rows, kept proportion = 0.9942, filtered \ensuremath{\approx} 116814*0.9942 \ensuremath{\approx} 116134.\\
llava\_instruct total 157710 rows, kept proportion =0.998 \ensuremath{\Rightarrow} \ensuremath{\approx} 157395.\\
ocrvqa total 165746 rows, kept proportion =0.9829 \ensuremath{\Rightarrow} \ensuremath{\approx} 162804.
\end{agentthought}

\noindent All five estimates are arithmetically correct. This reasoning was emergent:no tool or prompt instructed the agent to extrapolate from sample-level filter rates to full dataset sizes.

\subparagraph{Steps 7--8:Error recovery (lines 109--114).}
First \texttt{compute\_mix\_ratio} call passes 4 sizes for 5 datasets. The agent diagnoses and fixes in one step (line 111):

\begin{agentthought}[title={Agent Thought:self-correction}]
\small\ttfamily
The tool expects effective\_sizes count to match number of active datasets (including those with 0 rows?). It sees 5 active datasets (including captcha). We gave 4 sizes. Need to provide size for captcha (0).
\end{agentthought}

\begin{toolaction}[title={Tool Call (retry)}]
\small\ttfamily compute\_mix\_ratio(effective\_sizes="82540,0,116134,157395,162804", temperature=5, total\_samples=10000)\\
\textrm{$\to$ OK. Recipe: mode=auto\_temperature, $T{=}5$, total 9{,}998 samples.}
\end{toolaction}

\subparagraph{Steps 9--12:Mix, save, finetune, eval (lines 154--166).}
\mbox{}

\begin{toolaction}[title={Tool Calls}]
\scriptsize
\begin{tabular}{@{}rllr@{}}
\toprule
\textbf{Step} & \textbf{Call} & \textbf{Result} & \textbf{Time} \\
\midrule
9  & \texttt{vlm\_mix(counts="2296,0,2459,2615,2630")} & OK, 10{,}000 rows & 0.8s \\
10 & \texttt{save\_to\_disk(format="parquet")} & OK & 4.0s \\
11 & \texttt{submit\_finetune(dataset\_id=...)} & OK, job \texttt{ft-e95ca9d3} & 1{,}226s \\
12 & \texttt{submit\_eval(model\_id="ft-e95ca9d3")} & \textbf{ERR}, GPU OOM & 140s \\
\bottomrule
\end{tabular}\\[2pt]
\textrm{Scores recovered from cached VLMEvalKit outputs.}
\end{toolaction}

\noindent The agent passes the correct \texttt{job\_id} format (not a filesystem path) and omits training hyperparameters (fixed by task config). No hallucinated parameters across all 12 calls.

\subparagraph{Analysis.}
\begin{enumerate}
    \item \textbf{Cross-episode learning}: Reviews 4 prior episodes and formulates a hypothesis (that \texttt{quality\_filter} caused Ep~3's collapse). Whether the attribution is correct is unclear (see note above), but the resulting decision to skip text filtering is operationally sound.
    \item \textbf{Data-driven plan revision}: Drops captcha despite planning 10--15\%. Data overrides the prior plan.
    \item \textbf{Emergent quantitative reasoning}: Derives proportional effective-size estimates unprompted. All arithmetic correct.
    \item \textbf{Self-correction}: Diagnoses argument count mismatch and fixes in one step.
    \item \textbf{Efficiency}: 12 steps vs.\ 18--25 in Ep~1--4. Zero rollbacks, zero retried finetunes.
\end{enumerate}


%% ============================================================
%% NEXT STEPS
%% ============================================================

\meetingtasks{
    \item \todo{Stabilize eval/train subprocess lifecycle: GPU OOM from zombie processes caused 3/5 episode failures. Need process-group cleanup, GPU memory fence between finetune and eval, and automatic retry on OOM.}
    \item \todo{Run 10+ episode stress test: let the agent run overnight with seed memories and measure whether curation strategy improves from eval feedback across episodes.}
    \item \todo{Multi-agent systems literature review: 51 papers collected. Extract design decisions for sub-agent architecture (task decomposition, communication protocols, shared memory).}
}
